%%%%%%%%%%%%%%%%%
% CV tailored for Architecte IA (H/F) - Alpes-Maritimes
% Positioning: AI Architect / Senior Data Scientist in transition toward AI platform architecture
%%%%%%%%%%%%%%%%%

\documentclass[8pt,a4paper,ragged2e,withhyper]{../_shared/altacv}

\geometry{left=1.25cm,right=1.25cm,top=.5cm,bottom=1.cm,columnsep=1.2cm}

\usepackage{paracol}
\usepackage{fontawesome5}
\usepackage{hyperref}

\iftutex 
  \setmainfont{Roboto Slab}
  \setsansfont{Lato}
  \renewcommand{\familydefault}{\sfdefault}
\else
  \usepackage[rm]{roboto}
  \usepackage[defaultsans]{lato}
  \renewcommand{\familydefault}{\sfdefault}
\fi

\definecolor{SlateGrey}{HTML}{2E2E2E}
\definecolor{LightGrey}{HTML}{666666}
\definecolor{DarkPastelRed}{HTML}{8AC0D9}
\definecolor{PastelRed}{HTML}{00B0DA}
\definecolor{GoldenEarth}{HTML}{B4AD0C}

\colorlet{name}{black}
\colorlet{tagline}{PastelRed}
\colorlet{heading}{DarkPastelRed}
\colorlet{headingrule}{GoldenEarth}
\colorlet{subheading}{PastelRed}
\colorlet{accent}{PastelRed}
\colorlet{emphasis}{SlateGrey}
\colorlet{body}{LightGrey}

\definecolor{violet}{HTML}{5A2582}
\definecolor{rouge}{HTML}{F53B3B}
\definecolor{noir}{HTML}{00232C}
\definecolor{bleu}{HTML}{00B0DA}
\definecolor{rose}{HTML}{F831A0}
\definecolor{jaune}{HTML}{FFEC00}
\definecolor{vert}{HTML}{34AD6C}

\renewcommand{\namefont}{\Huge\rmfamily\bfseries}
\renewcommand{\personalinfofont}{\footnotesize}
\renewcommand{\cvsectionfont}{\LARGE\rmfamily\bfseries}
\renewcommand{\cvsubsectionfont}{\large\bfseries}
\renewcommand{\cvItemMarker}{{\small\textbullet}}
\renewcommand{\cvRatingMarker}{\faCircle}

\input{../_shared/pubs-num.tex}
\addbibresource{../_shared/sample.bib}

\begin{document}

\name{BOILEAU Guillaume}
\tagline{Architecte IA / Senior Data Scientist | Python, ML, LLMs, MLOps foundations, systèmes complexes}

\photoL{2.8cm}{../_shared/moi}

\personalinfo{%
  \email{guillaume.boileaupro@gmail.com}
  \phone{+33 6 59 94 85 84}
  \location{Alpes-Maritimes, FRANCE}
  \linkedin{guillaume-boileau-891b81265}
  \researchgate{Guillaume-Boileau-2}
  \orcid{0000-0002-3576-6968}
}

\makecvheader
\vspace{-0.3cm}

\AtBeginEnvironment{itemize}{\small}
\columnratio{0.64}

\begin{paracol}{2}

\cvsection{Experience}

\cvevent{\textbf{Software Engineer -- AI-assisted development, integration \& operational validation}}{VEV Platform Services France}{2025 -- 2026}{Valbonne, France}

\textbf{Context:} Industrial software environment for an EV charging CPMS platform connected to operational charging stations, backend services, field equipment and hardware/software interfaces.

\textbf{Objective:} Contribute to software delivery, integration and operational reliability by combining development, data-flow analysis, incident investigation, technical documentation and AI-assisted engineering practices.

\begin{itemize}
  \item \textbf{Software delivery:} Contributed to TypeScript, Angular and Node.js components integrated into an operational software platform.
  \item \textbf{AI-assisted workflows:} Used AI tools for code review, debugging support, documentation, incremental problem solving and productivity improvement.
  \item \textbf{System integration:} Worked on interfaces between software services, operational data, connected infrastructure and field constraints.
  \item \textbf{Operational data flows:} Analysed FTP-based exchanges, interrupted flows, inconsistent data and service-impacting behaviours.
  \item \textbf{Validation:} Checked consistency between application data, expected behaviours and operational constraints.
  \item \textbf{Monitoring and support:} Investigated incidents, documented technical findings and contributed to service reliability improvements.
  \item \textbf{Agile coordination:} Used Zenhub for task tracking, backlog follow-up, prioritisation and coordination in an iterative environment.
\end{itemize}

\divider

\cvevent{\textbf{Senior R\&D Engineer -- Scientific AI, model validation \& Python platform development}}{CNES}{2023 -- 2025}{Nice, France}

\textbf{Context:} R\&D and engineering activities for the LISA ESA/NASA space mission, involving simulation, scientific Python, statistical modelling, signal processing, model integration and validation workflows.

\textbf{Objective:} Design and validate reproducible Python workflows able to integrate heterogeneous models, compare outputs, assess weak signals under noise and produce traceable results for collaborative technical review.

\begin{itemize}
  \item \textbf{Python platform:} Developed CATpyLISA, a Python platform for model integration, comparison, validation and technical review.
  \textcolor{DarkPastelRed}{\href{https://gitlab.oca.eu/gboileau/lisa_l3_tools}{\bf GitLab: CATpyLISA}}
  \item \textbf{Data science workflows:} Structured datasets, simulations, model outputs, validation logic and reproducible analysis pipelines.
  \item \textbf{Statistical modelling:} Worked on low signal-to-noise physical backgrounds, uncertainty analysis and model comparison.
  \item \textbf{Machine learning mindset:} Applied data-driven reasoning, feature interpretation, validation metrics and reproducibility principles to complex scientific data.
  \item \textbf{Model validation:} Defined consistency checks, comparison rules, acceptance criteria and validation procedures for heterogeneous contributions.
  \item \textbf{Technical coordination:} Interfaced with scientific, software and institutional contributors in an international space-mission environment.
  \item \textbf{Documentation and knowledge sharing:} Used Confluence-style documentation logic for traceability, review preparation and collaborative knowledge transfer.
  \item \textbf{Decision support:} Produced technical syntheses, identified discrepancies and formalised recommendations for model and workflow improvements.
\end{itemize}

\divider

\cvevent{\textbf{Data Scientist -- Statistical modelling, noise analysis \& ML-oriented workflows}}{Universiteit Antwerpen, Belgium}{2021 -- 2023}{Antwerp, Belgium}

\textbf{Context:} Data science and performance studies for next-generation precision instruments, combining environmental sensor data, statistical modelling, weak-signal detection, noise reduction and validation of complex systems.

\textbf{Objective:} Develop robust and reproducible workflows to identify, characterise and reduce noise sources affecting measurement quality and system performance.

\begin{itemize}
  \item \textbf{Data science:} Developed Python workflows for signal analysis, statistical modelling, performance assessment and validation.
  \item \textbf{Bayesian modelling:} Built MCMC-based pipelines to analyse different detector configurations and noise scenarios.
  \item \textbf{Sensor data:} Analysed seismometer and environmental data, propagation effects, spatial correlations and noise impacts.
  \item \textbf{Machine learning:} Structured approaches for non-stationary noise reduction using witness channels, machine learning and deep learning.
  \item \textbf{Performance metrics:} Used figures of merit, sensitivity studies and validation checks to compare configurations and quantify robustness.
  \item \textbf{Technical reporting:} Produced traceable results, technical syntheses and recommendations for international stakeholders.
\end{itemize}

\divider

\cvevent{\textbf{Junior R\&D Engineer -- Signal processing, simulation \& scientific computing}}{ARTEMIS, Observatoire de la Côte d’Azur}{2018 -- 2021}{Nice, France}

\textbf{Context:} Scientific computing and signal-processing activities for gravitational-wave mission preparation, involving simulation, weak-signal separation, diagnostic analysis and uncertainty studies.

\textbf{Objective:} Develop robust numerical methods for weak-signal analysis, simulation, uncertainty assessment and interpretation of complex datasets.

\begin{itemize}
  \item \textbf{Signal processing:} Developed methods to analyse and separate overlapping low-SNR signals.
  \item \textbf{Simulation:} Built numerical approaches for complex signal environments and large-scale scientific datasets.
  \item \textbf{Diagnostics:} Investigated noise sources and correlations across sensors and simulated observables.
  \item \textbf{Scientific Python:} Developed reproducible Python-based analysis, modelling and validation workflows.
  \item \textbf{Documentation:} Formalised methods, assumptions, results and limitations for collaborative review.
\end{itemize}

\switchcolumn

\cvsection{Profile}

\textbf{PhD physicist and senior data scientist with experience in Python, statistical modelling, machine learning, signal processing, model validation and software workflows for complex systems.} Strong background in scientific AI, reproducible pipelines, low-SNR data, technical coordination and knowledge transfer.

Positioning for this role: AI architect / senior data scientist able to structure AI use cases from ideation to validated POCs, support industrialisation with MLOps foundations, and help teams adopt LLM, predictive analytics and data-driven workflows.

\begin{itemize}
  \item Strong quantitative background: statistics, modelling, signal processing, uncertainty analysis and validation.
  \item Python development for reproducible data-processing, simulation and model-validation pipelines.
  \item Experience with ML-oriented workflows, deep learning concepts, LLM/agentic AI tools, chatbots and MCP-based architectures.
  \item Able to bridge technical architecture, data science, software delivery, documentation and stakeholder acculturation.
  \item Fast ramp-up profile for Azure ML, Databricks, MLflow and enterprise LLMOps environments.
\end{itemize}

\cvsection{Languages}

\cvskill{English}{4}
\divider
\cvskill{Spanish}{2}
\divider

\medskip

\cvsection{Education}

\cvevent{PhD in Physics}{Université Côte d’Azur}{2018 -- 2021}{Nice, France}

\divider

\cvevent{Selective Graduate Program in Fundamental Physics (Magistère)}{Université Paris-Saclay}{2016 -- 2018}{Orsay, France}

\divider

\cvevent{MSc in Astronomy and Astrophysics}{Observatoire de Paris / Université Paris-Saclay}{2017 -- 2018}{Paris, France}

\divider

\cvevent{BSc in Fundamental Physics}{Université Paris-Saclay}{2015 -- 2016}{Orsay, France}

\cvsection{Professional Training}

\cvevent{Machine Learning in Python with scikit-learn}{FUN MOOC -- Inria}{2026}{France Université Numérique}

\small{
Predictive modelling, supervised learning, model evaluation, cross-validation, metrics, pipelines and practical machine learning workflows with Python and scikit-learn.
}

\divider

\cvevent{Supervised Machine Learning: Regression \& Classification}{DeepLearning.AI -- Andrew Ng}{2020}{Coursera}

\divider

\cvevent{Recherche reproductible : principes méthodologiques pour une science transparente}{FUN MOOC -- Inria}{2025}{France Université Numérique}

\divider

\cvevent{Continuous Professional Training -- Software, Data, AI and Systems}{OpenClassrooms}{2024 -- 2026}{}

\small{
Python, SQL, Machine Learning, AI, Linux, JavaScript, TypeScript, Angular, Node.js, Django, REST APIs, C/C++, web development, algorithms and software engineering fundamentals.
}

\end{paracol}

\cvsection{Projects}

\cvevent{CATpyLISA -- Python platform for model integration, comparison and validation}{ESA/NASA LISA space mission}{2023 -- 2025}{}

Python platform dedicated to heterogeneous model integration, data structuring, comparison of outputs, validation logic and collaborative technical review for a large-scale international scientific mission.

\textbf{Role:} Python development, pipeline design, data structuring, model validation, consistency checks, documentation, technical coordination and contribution consolidation.

\textbf{Skills:} Python, data workflows, model validation, signal processing, statistical modelling, simulation, documentation, GitLab, reproducibility, international collaboration.

\divider

\cvevent{AI-assisted engineering, chatbots \& agentic workflows}{Software productivity / LLM tooling / MCP architectures}{2024 -- 2026}{}

Use and exploration of AI-assisted development workflows for code review, debugging, documentation, technical synthesis, incremental reasoning, chatbot logic and tool-oriented agent architectures.

\textbf{Role:} prompting, code review support, architecture reasoning, workflow design, documentation and evaluation of AI-assisted productivity patterns.

\textbf{Skills:} LLMs, agents, MCP, chatbot architecture, prompting, Python, software engineering, documentation, productivity tooling.

\divider

\cvevent{Noise reduction and predictive modelling workflows}{Machine learning / sensor data / precision instrumentation}{2021 -- 2023}{}

Development of numerical and statistical workflows for weak-signal detection, sensor-data analysis, non-stationary noise reduction and performance comparison across system configurations.

\textbf{Role:} data analysis, statistical modelling, MCMC pipelines, sensor-data interpretation, machine-learning-oriented noise mitigation and performance validation.

\textbf{Skills:} Python, statistics, MCMC, signal processing, ML concepts, deep learning concepts, sensor data, validation, technical reporting.

\cvsection{Strengths}

\vspace{-.3cm}

\cvsubsection{AI architecture, LLMs \& acculturation}

{\small
\cvtag{AI Architecture -- ramp-up}
\cvtag{AI Strategy}
\cvtag{LLMs}
\cvtag{Agentic AI}
\cvtag{MCP}
\cvtag{Chatbots}
\cvtag{Prompt Engineering}
\cvtag{Use-case Framing}
\cvtag{POC Coordination}
\cvtag{AI Acculturation}
\cvtag{Training Workshops}
\cvtag{Technical Vulgarisation}
\cvtag{Change Management}
\cvtag{Best Practices}
}

\divider\smallskip
\vspace{-.3cm}

\cvsubsection{Data science, machine learning \& modelling}

{\small
\cvtag{Data Science}
\cvtag{Machine Learning}
\cvtag{Predictive Analytics}
\cvtag{Statistical Modelling}
\cvtag{Bayesian Inference}
\cvtag{MCMC}
\cvtag{Regression}
\cvtag{Classification}
\cvtag{Model Evaluation}
\cvtag{Cross-Validation}
\cvtag{Metrics}
\cvtag{Feature Engineering}
\cvtag{Uncertainty Analysis}
\cvtag{Computer Vision -- ramp-up}
\cvtag{NLP -- ramp-up}
}

\divider\smallskip
\vspace{-.3cm}

\cvsubsection{MLOps / LLMOps foundations}

{\small
\cvtag{MLOps Foundations}
\cvtag{LLMOps -- ramp-up}
\cvtag{CI/CD}
\cvtag{Versioning}
\cvtag{Monitoring}
\cvtag{Reproducible Workflows}
\cvtag{Model Validation}
\cvtag{Data Quality}
\cvtag{Traceability}
\cvtag{Docker}
\cvtag{Git / GitLab}
\cvtag{Linux}
\cvtag{Security \& Compliance -- awareness}
\cvtag{Azure ML -- ramp-up}
\cvtag{Databricks -- ramp-up}
\cvtag{MLflow -- ramp-up}
}

\divider\smallskip
\vspace{-.3cm}

\cvsubsection{Python, software \& data engineering}

{\small
\cvtag{Python}
\cvtag{NumPy}
\cvtag{SciPy}
\cvtag{pandas}
\cvtag{scikit-learn}
\cvtag{Matplotlib}
\cvtag{Jupyter}
\cvtag{SQL}
\cvtag{Bash}
\cvtag{REST API}
\cvtag{TypeScript}
\cvtag{Angular}
\cvtag{Node.js}
\cvtag{C}
\cvtag{C++}
}

\divider\smallskip
\vspace{-.3cm}

\cvsubsection{Systems, validation \& industrialisation}

{\small
\cvtag{Complex Systems}
\cvtag{System Integration}
\cvtag{Model Integration}
\cvtag{Validation Strategy}
\cvtag{IVVQ}
\cvtag{Anomaly Analysis}
\cvtag{Root Cause Analysis}
\cvtag{Technical Reporting}
\cvtag{Decision Support}
\cvtag{Requirements Awareness}
\cvtag{Risk Analysis}
\cvtag{Industrial Software}
\cvtag{Production Support}
\cvtag{Operational Monitoring}
}

\divider\smallskip
\vspace{-.3cm}

\cvsubsection{Tools, collaboration \& governance}

{\small
\cvtag{Confluence}
\cvtag{Jira}
\cvtag{Zenhub}
\cvtag{OpenSearch}
\cvtag{Sentry}
\cvtag{Power BI}
\cvtag{Agile}
\cvtag{Kanban}
\cvtag{Stakeholder Coordination}
\cvtag{Documentation}
\cvtag{Knowledge Sharing}
\cvtag{International Environment}
\cvtag{Technical English}
}

\end{document}
